\documentclass[11pt]{article}

\usepackage{amsmath, amssymb, amsthm}
\usepackage{graphicx}
\usepackage{hyperref}
\usepackage{geometry}
\geometry{margin=1in}

\title{Representation Entropy as a Unifying Framework for\\
Classical and Quantum Efficiency in Molecular Simulation}

\author{Fábio Novaes\footnote{Federal Rural University of Pernambuco, fabio.nsantos@gmail.com}}

\date{\today}

\begin{document}

\maketitle

\begin{abstract}
% TODO:
% - 3–5 frases
% - Problema: vantagem quântica em química é incerta
% - Ideia: entropia de representação (Fourier + Pauli)
% - Resultado: diagrama de fases
% - Validação: experimentos simples (H2, LiH)
\end{abstract}

% ============================================================
\section{Introduction}

% TODO:
% - Motivação: química quântica como alvo de quantum advantage
% - Evidências conflitantes:
%   - Chan et al.: sem vantagem exponencial
%   - Experimentos recentes: sucesso híbrido
% - Problema: falta de critério unificado

% TODO:
% - Contribuição do paper:
%   - Definir entropia de representação
%   - Unificar Fourier e Pauli
%   - Explicar VQE, métodos clássicos e híbridos
%   - Propor phase diagram

% TODO:
% - Contributions (bullet points):
%   1. Definição formal de H_rep
%   2. Relação com trainability e simulabilidade
%   3. Phase diagram
%   4. Validação numérica

% ============================================================
\section{Background}

\subsection{Variational Quantum Algorithms}

% TODO:
% - Definir VQE
\begin{equation}
E(\theta) = \langle \psi(\theta) | H | \psi(\theta) \rangle
\end{equation}

% TODO:
% - Problemas:
%   - barren plateaus
%   - custo de medição

\subsection{Classical Methods}

% TODO:
% - CC, DMRG, HCI
% - Ideia central: compressibilidade

\subsection{Recent Results}

% TODO:
% - Nemkov (Fourier)
% - Okumura (barren plateau = Fourier)
% - Angrisani (Pauli truncation)
% - Chan (no exponential advantage)
% - Quantum-centric experiment

% ============================================================
\section{Fourier and Pauli Representations}

\subsection{Fourier Expansion}

\begin{equation}
E(\theta) = \sum_{k} c_k e^{i k \cdot \theta}
\end{equation}

% TODO:
% - explicar origem (geradores do circuito)
% - discutir sparsidade

\subsection{Pauli Expansion}

\begin{equation}
U^\dagger O U = \sum_{\gamma} \Phi_\gamma P_\gamma
\end{equation}

% TODO:
% - explicar Heisenberg picture
% - conectar com Angrisani

\subsection{Relation Between Representations}

% TODO:
% - mostrar que c_k = combinação linear de Phi_gamma
% - discutir mudança de base

% ============================================================
\section{Representation Entropy}

\subsection{Fourier Entropy}

\begin{equation}
p(k) = \frac{|c_k|^2}{\sum_{k'} |c_{k'}|^2}
\end{equation}

\begin{equation}
H_F = -\sum_k p(k)\log p(k)
\end{equation}

% TODO:
% - interpretação: complexidade espectral

\subsection{Pauli Entropy}

\begin{equation}
p(\gamma) = \frac{|\Phi_\gamma|^2}{\sum_{\gamma'} |\Phi_{\gamma'}|^2}
\end{equation}

\begin{equation}
H_P = -\sum_\gamma p(\gamma)\log p(\gamma)
\end{equation}

% TODO:
% - interpretação: spreading do operador

\subsection{Unified Representation Entropy}

% TODO:
% - discutir equivalência aproximada
% - definir H_rep

% ============================================================
\section{Phase Diagram of Variational Quantum Algorithms}

\subsection{Axes of the Diagram}

% TODO:
% - eixo x: H_F (trainability)
% - eixo y: H_P (simulabilidade)

\subsection{Regimes}

\paragraph{Region I: Classically Tractable}
% TODO:
% H_F baixo, H_P baixo

\paragraph{Region II: Variationally Useful}
% TODO:
% H_F moderado, H_P moderado

\paragraph{Region III: Barren Plateau}
% TODO:
% H_F máximo

\paragraph{Region IV: Potential Quantum Advantage}
% TODO:
% H_P alto, H_F moderado

\paragraph{Region V: Intractable}
% TODO:
% ambos altos

\subsection{Scaling Relations}

\begin{equation}
T \sim e^{-H_F}
\end{equation}

\begin{equation}
C \sim e^{-H_P}
\end{equation}

% TODO:
% interpretar essas relações

% ============================================================
\section{Connection to Existing Results}

\subsection{Barren Plateaus}

% TODO:
% usar Okumura:
\begin{equation}
\sum_k |c_k|^2 \sim 2^{-n}
\end{equation}

% TODO:
% mostrar H_F máximo

\subsection{Classical Simulability}

% TODO:
% Chan + Angrisani
% H_P baixo → truncável

\subsection{Hybrid Quantum Methods}

% TODO:
% quantum-centric experiment
% subspace ≈ suporte espectral

% ============================================================
\section{Numerical Experiments}

\subsection{Setup}

% TODO:
% - sistemas: H2, LiH
% - método: Qiskit + PySCF
% - ansatz: UCC / hardware-efficient

\subsection{Estimating Fourier Coefficients}

% TODO:
% - amostrar E(theta)
% - ajustar série de Fourier

\subsection{Estimating Representation Entropy}

% TODO:
% - calcular H_F
% - estimar H_P (truncado)

\subsection{Results}

% TODO:
% - plot H_F vs distância
% - plot erro VQE vs H_F
% - discutir correlação

% \begin{figure}
% \centering
% \includegraphics[width=0.6\textwidth]{phase_diagram.png}
% \caption{Phase diagram of VQAs in the $(H_F, H_P)$ plane.}
% \end{figure}

% ============================================================
\section{Discussion}

% TODO:
% - interpretação física
% - por que química é compressível
% - implicações para design de ansatz

% TODO:
% - limitações:
%   - cálculo de entropia
%   - sistemas pequenos

% TODO:
% - implicações:
%   - critério de vantagem quântica

% ============================================================
\section{Conclusion}

% TODO:
% - resumo dos resultados
% - frase forte:
%   complexidade = entropia de representação

% TODO:
% - perspectivas futuras:
%   - dinâmica
%   - novos ansatz
%   - sistemas fortemente correlacionados

% ============================================================
\appendix

\section{Derivations}

% TODO:
% - Fourier ↔ Pauli
% - detalhes matemáticos

\section{Additional Data}

% TODO:
% - figuras extras
% - tabelas

% ============================================================

% \bibliographystyle{plain}
% \bibliography{refs}

\end{document}